\emph{Exact fibre interval of the diagonal quadratic statistic and scalar non-recovery from the Gram spectrum.} Within the named constant-variance gaussian edge-block subclass of the second preceding paragraph and the column-profile edge-block law of the first preceding paragraph at every strictly positive non-constant variance profile, fix \(D\ge3\) and use the coordinates of those paragraphs: \(\tau=\operatorname{Tr}G\), \(p=\lambda/\tau\), \(d_{i}=G_{ii}/\tau\), and the single scalar \(P_{2}=\sum_{i}d_{i}^{2}\), the \(D=2\) instance of which is the statistic \(\hat\rho_{2}\) of the third preceding paragraph. Two elementary identities hold for every Hermitian \(G\): writing \(\sum_{j\ne k}\lvert G_{jk}\rvert^{2}\) for the sum of squared moduli of the off-diagonal entries, one has \(\tau^{2}P_{2}=\sum_{j}\lambda_{j}^{2}-\sum_{j\ne k}\lvert G_{jk}\rvert^{2}\), so \(P_{2}\le\sum_{j}p_{j}^{2}\) with equality exactly for diagonal \(G\); and \(P_{2}\ge1/D\) by Cauchy's inequality \(\sum_{i}d_{i}^{2}\ge(\sum_{i}d_{i})^{2}/D\), with equality exactly for the uniform diagonal \(d_{i}=1/D\). The value set of \(P_{2}\) on the unitary orbit of \(G\) is then exactly the closed interval from \(1/D\) to \(\sum_{j}p_{j}^{2}\): the unitary group is connected, so the orbit is connected and the image is an interval; the upper endpoint is attained at the diagonal point of the orbit; and the lower endpoint is attained at the Fourier point \(F^{\dagger}\Lambda F\), where \(\Lambda=\operatorname{diag}(\lambda)\) and \(F\) is the character table of the cyclic group of order \(D\) divided by \(\sqrt{D}\), whose squared moduli are all \(1/D\), so its diagonal is uniform. The interval is non-degenerate exactly when the spectrum is non-scalar, and at \(D=2\) it is exactly the fibre interval from \(1/2\) to \((\lambda_{1}^{2}+\lambda_{2}^{2})/\tau^{2}\) of the third preceding paragraph. This pointwise interval is all the measure-theoretic step requires. Under either named law the eigenvalue-eigenbasis disintegration of the second and first preceding paragraphs applies: conditional on the spectrum the eigenbasis law is the Haar fibre reference in the constant-variance class and, under every strictly positive profile, a density bounded above and below by positive constants against that reference in the profile class; in both cases the conditional measure of the fibre has full support on it, charging every nonempty orbit-neighborhood, and the eigenvalue-collision set still carries no mass, the joint eigenvalue density being absolutely continuous on the ordered chamber with the Vandermonde weight recorded in the first preceding paragraph. If a map \(F_{0}\) of the Gram-spectrum coordinate recovered \(P_{2}\) almost everywhere, then for almost every spectrum the restriction of \(P_{2}\) to the fibre orbit would equal the constant \(F_{0}(p)\) off a set of fibre measure zero; that restriction is continuous, so it would equal the constant everywhere, since its deviation from the constant, if nonzero at one point, is nonzero on a nonempty orbit-neighborhood of positive fibre measure; and at almost every spectrum the interval is non-degenerate, a contradiction. Hence \(\mathbb{P}(P_{2}=F_{0}(p))=0\) for every map \(F_{0}\), for the law of every datum in either named class at every \(D\ge3\); combined with the \(D=2\) nonexistence of the third preceding paragraph and the \(D=2\) instance of the tuple nonexistence of the first preceding paragraph, this settles the single quadratic statistic at every dimension \(D\ge2\) in both named classes, the scalar being exactly as invisible to the Gram spectrum as the full diagonal tuple of the second and first preceding paragraphs. The contrast with the fibre dimension is exact and elementary: adding this scalar to the spectrum saturates the fibre only at \(D=2\), where the diagonal fibre is the interval of \(G_{11}\) of the third preceding paragraph and \(\{d_{1},d_{2}\}=\{(1+\sqrt{2P_{2}-1})/2,\,(1-\sqrt{2P_{2}-1})/2\}\) recovers the diagonal multiset from \((p,P_{2})\); at \(D\ge3\) the scalar does not saturate the fibre even modulo permutation, because, for every \(D\ge3\), intersecting the two-angle rotation patch of the second preceding paragraph - at \(D>3\) the same two-angle construction taken inside a three-dimensional eigen-coordinate block with the remaining diagonal entries held fixed, the two-angle Jacobian computation being block-local - with a level set of \(P_{2}\) in the interior of its image range under the patch yields an uncountable family of distinct diagonals, only finitely many of which are related by a permutation of the \(D\) entries, so the power-sum and \(\rho\)-family tuple negatives of the second and first preceding paragraphs remain strictly stronger statements than this scalar one. Not claimed here: the exact conditional density of \(P_{2}\) given the spectrum beyond the displayed \(D=2\) multiset formula, the joint law of \(P_{2}\) with the remaining power sums, and any identification of fibre diagonal images with the full Schur--Horn permutahedron, which remains outside the proved content of the preceding paragraphs and of this one. These statements are exact linear-algebraic and measure-theoretic identifiability statements on the support of the named constant-variance subclass and of the named column-profile law at every strictly positive non-constant profile and every \(D\ge3\), they bound nonexistence between the named scalar and spectrum coordinates only, they evaluate no matrix integral of the core action, they supply no state and no selection mechanism, they do not fix the state that sets \(W\), \(g\), \(s\) or the frame, they leave the missing steps (i)-(iv) of the Missing step paragraph unchanged, and they close nothing in Problem~G.
